\section{Conclusion}
\label{sec:conclusion}
Host-based detection rules are code with kernel visibility, and today they
run as trusted code. \sys runs them in WebAssembly instead and this paper
measures what that costs and what it buys. The boundary itself is cheap:
\eTwoEngineOverheadUs{}~\si{\micro\second} per call empty,
\eTwoMlRawUs{}~\si{\micro\second} for a real rule on a binary-encoded event,
a sustained \eThreeAchievedBinFiftyKRf{} syscall events per second through
the whole pipeline, and an application cost under nginx, redis and
PostgreSQL in the same range as Falco's. It holds: ten adversarial modules
in eleven attempts were contained by memory, fuel and deadline limits with
their neighbours unaffected. And it makes a claim checkable that is
usually taken on faith: the detector the system runs is, to
\eTwoMaxDiffNumpyWasm{} for the MLP and \eOneWasmSklearnMaxDiff{} for the
forest, the detector that was evaluated, which is how we can report that
\sys's accuracy on ADFA-LD is an AUPRC of \eOneAuprcRfWasm{} with the forest
and \eOneAuprcMlp{} with the MLP it first shipped with, and how the better
model got into the module in a day. The remaining cost is in the rules, not
the sandbox, and we have named the two levers that move it.
